\section{Weil's distribution and the damped criterion}

\subsection{The deconvolved object exists: it is Weil's distribution}\label{dc:sec:weil}

With $\hat f(t)=\int f(x)e^{itx}dx$ the Fourier transform of $P_a$ is $e^{-a|x|}$, so deconvolution by $P_a$ multiplies Fourier transforms by $e^{a|x|}$. Let $\mathcal B$ be the space of even functions $\varphi(\theta)=\frac1{2\pi}\int\hat\varphi(x)e^{-i\theta x}dx$ with $\hat\varphi$ of compact support, Lipschitz, and with $\hat\varphi'$ of bounded variation. Then $\varphi$ is entire of exponential type and $\varphi(\theta+is)=O(\theta^{-2})$ uniformly for $|s|\le1$, since $\hat\varphi(x)e^{sx}$ has the same properties. For $a\ge0$ let $P_a^{-1}\varphi\in\mathcal B$ have Fourier transform $\hat\varphi(x)e^{a|x|}$; the corner of $e^{a|x|}$ at $0$ keeps it in $\mathcal B$.

\begin{theorem}[The deconvolved functional]\label{dc:thm:weil}
Let $\varepsilon_0$ be to the right of the zeros. For $\varphi\in\mathcal B$ put $D(\varphi)=\int\rho_{\varepsilon_0}(\theta)\,(P_{\varepsilon_0}^{-1}\varphi)(\theta)\,d\theta$. Then $D$ does not depend on $\varepsilon_0$, and with $g(x)=\hat\varphi(x)/2\pi$, so that $\varphi(\theta)=\int g(x)e^{i\theta x}dx=\hat g(\theta)$,
\[
\begin{aligned}
D(\varphi)=\sum_\rho\varphi(\gamma_\rho)={}&\hat g\big(\tfrac i2\big)+\hat g\big(-\tfrac i2\big)-(\gamma_E+\log\pi)\,g(0)\\
&+\int_0^\infty\frac{2\big(e^{-2x}g(0)-e^{-x/2}g(x)\big)}{1-e^{-2x}}\,dx\\
&-2\sum_{n\ge2}\frac{\Lambda(n)}{\sqrt n}\,g(\log n),
\end{aligned}
\]
where the sum over $\rho$ runs over all zeros with multiplicity, in conjugate pairs, and converges absolutely. Through the $x$ form of its archimedean term, the Thorin functional $U$ of Section~\ref{sec:ch14:functional} extends to $\mathcal B$, and then $D(\varphi)=2\langle U,\varphi(\sqrt{\cdot})\rangle$.
\end{theorem}

\begin{proof}
By the expansion of $\rho_{\varepsilon_0}$ over the zeros (Proposition~\ref{prop:struct}, summed over all zeros $\rho=\beta+i\gamma$, with $\gamma$ of both signs), $D(\varphi)=\sum_\rho\int P_{c_\rho}(\theta-\gamma)\,(P_{\varepsilon_0}^{-1}\varphi)(\theta)\,d\theta$ with $c_\rho=\varepsilon_0-(\beta-\frac12)>0$; the interchange is justified because $\rho_{\varepsilon_0}=O(\log(2+|\theta|))$ and $P_{\varepsilon_0}^{-1}\varphi=O(\theta^{-2})$. By Parseval each term is $\frac1{2\pi}\int\hat\varphi(x)e^{(\varepsilon_0-c_\rho)|x|}e^{-i\gamma x}\,dx$. A zero $\rho$ off the line and its partner $1-\bar\rho$ have the same $\gamma$ and opposite $\beta-\frac12=\pm\delta$, so the pair gives $\frac1{2\pi}\int\hat\varphi(x)e^{-i\gamma x}\,2\cosh(\delta x)\,dx=\varphi(\gamma-i\delta)+\varphi(\gamma+i\delta)$, because $e^{\delta|x|}+e^{-\delta|x|}=e^{\delta x}+e^{-\delta x}$; a zero on the line is its own partner and gives the single term $\varphi(\gamma)$. This is $\sum_\rho\varphi(\gamma_\rho)$, absolutely convergent by the uniform $O(\theta^{-2})$ bound, and it does not involve $\varepsilon_0$. The second expression is the explicit formula of Weil~\cite{Weil52}, \cite[Theorem 5.12]{IwaniecKowalski04}, with the archimedean term written in the variable $x$ by $\re\psi(\frac14+\frac{it}2)=-\gamma_E+\int_0^\infty\frac{e^{-u}-e^{-u/4}\cos(tu/2)}{1-e^{-u}}\,du$; this integral converges because $g$ is Lipschitz at $0$, and the formula extends from the usual class to $\mathcal B$ by approximation: mollifying $g$ gives functions of the usual class supported in a fixed interval, with uniformly bounded Lipschitz constants, and both sides are continuous along such sequences, the zero side by the uniform $O(\theta^{-2})$ bound. The comparison with~\eqref{eq:ch14:U}, whose functional sums once over each pair $\{\rho,1-\rho\}$, gives the factor $2$.
\end{proof}

For $g$ even, Lipschitz and of compact support, with $g'$ of bounded variation, so that $\hat g\in\mathcal B$, write $\langle D,g\rangle=D(\hat g)$, and for $g$ not even apply it to the even part of $g$. By the theorem $\langle D,g\rangle=\sum_\rho\hat g(\gamma_\rho)$, which is the action on $g$ of Weil's distribution $D(x)=\sum_\rho e^{(\rho-\frac12)x}$ of Proposition~\ref{rz:thm:weil}.

The distribution $P_{\varepsilon_0}^{-1}\rho_{\varepsilon_0}$ on $\mathcal B$, which exists with no hypothesis, is a positive measure exactly when RH holds: under RH it is $\sum_\gamma\delta_\gamma$, and if $D(\varphi)=\int\varphi\,d\mu$ on $\mathcal B$ with $\mu\ge0$, then $\int\rho_{\varepsilon_0}\psi=D(P_{\varepsilon_0}*\psi)=\int(P_{\varepsilon_0}*\psi)\,d\mu$ for every $\psi\in\mathcal B$, which gives $\int(1+\theta^2)^{-1}d\mu<\infty$ (take $\psi\ge0$ with $P_{\varepsilon_0}*\psi\ge c(1+\theta^2)^{-1}$) and $\rho_{\varepsilon_0}=P_{\varepsilon_0}*\mu$, so that RH follows from Theorem~\ref{tf:thm:deconv}. On the functions $\varphi=|\hat h|^2$ the statement $D(\varphi)\ge0$ is Weil's criterion~\cite{Weil52}, in the form the book uses in Section~\ref{sec:ch10:positivity}. We checked the second expression against $2\sum_{\gamma>0}\hat g(\gamma)$ for the Gaussian $g(x)=e^{-(x/c)^2}$ (\texttt{code/dc\_weil.py}): at $c=\frac12$ both sides are $6.69239857682\times10^{-6}$ and they differ by less than $10^{-69}$ at seventy digits of working precision.

\subsection{The damped Weil criterion}\label{dc:sec:damped}

Partial deconvolution, to depth $\varepsilon$ instead of to the line, corresponds to damping. For bounded $h$ of bounded variation and compact support put $\tilde h(x)=\overline{h(-x)}$, so that $h*\tilde h$ is Lipschitz and Hermitian, and
\[
Q_\varepsilon(h)=\big\langle D,\,e^{-\varepsilon|x|}\,(h*\tilde h)\big\rangle .
\]
The Fourier transform of $e^{-\varepsilon|x|}$ is $2\pi P_\varepsilon$ and that of $h*\tilde h$ is $|\hat h|^2$, so
\[
\widehat{e^{-\varepsilon|\cdot|}(h*\tilde h)}=P_\varepsilon*|\hat h|^2 ,
\]
a nonnegative function. Real $h$ suffice in what follows, since $Q_\varepsilon(h)=Q_\varepsilon(\re h)+Q_\varepsilon(\im h)$.

\begin{theorem}[Damped Weil criterion]\label{dc:thm:damped}
Let $0\le\varepsilon\le\frac12$. Then $Q_\varepsilon(h)\ge0$ for every $h\in C_c^\infty(\R)$ if and only if no nontrivial zero of $\zeta$ has $\beta>\frac12+\varepsilon$. When $\varepsilon$ is to the right of the zeros,
\begin{equation}\label{dc:eq:Qw}
Q_\varepsilon(h)=\int_{\R}|\hat h(\theta)|^2\,\rho_\varepsilon(\theta)\,d\theta
\end{equation}
for every bounded $h$ of bounded variation and compact support. In particular: at $\varepsilon=0$ this is Weil's criterion for RH; at $\varepsilon=\frac12$ it holds with no hypothesis; and the zero free half plane $\re s>\frac78$ announced in October 2026~\cite{OAI26a}, which has not been refereed, is equivalent to $Q_{3/8}\ge0$.
\end{theorem}

\begin{proof}
Formula~\eqref{dc:eq:Qw} follows from Theorem~\ref{dc:thm:weil}. Let $h$ be real, bounded, of bounded variation and of compact support. Then $F=h*\tilde h$ is Lipschitz and of compact support, and its derivative $h*d\tilde h$, the convolution of a bounded function of bounded variation with a finite measure, is of bounded variation; so $e^{-\varepsilon|x|}F$ is Lipschitz, of compact support, with a derivative of bounded variation, and $\varphi=P_\varepsilon*|\hat h|^2$, the Fourier transform of $e^{-\varepsilon|x|}F$, and $|\hat h|^2$, that of $F$, lie in $\mathcal B$. As $P_\varepsilon^{-1}\varphi=|\hat h|^2$, the first identity of Theorem~\ref{dc:thm:weil} with $\varepsilon_0=\varepsilon$ gives $Q_\varepsilon(h)=D(\varphi)=\int|\hat h|^2\rho_\varepsilon\,d\theta$; for complex $h$ add the two parts, the cross terms of $|\hat h|^2$ being odd and $\rho_\varepsilon$ even. For bounded piecewise $C^1$ functions $h$ this is Proposition~\ref{yo:prop:mult} of Chapter~\ref{chap:yoshida} at $\alpha=\frac12+\varepsilon$, where the form $\langle\varphi,\varphi\rangle_\alpha$ is $Q_\varepsilon(\varphi)$. The \emph{if} direction is Theorem~\ref{yo:thm:damped}(i): if no zero has $\beta>\frac12+\varepsilon$, then either $\varepsilon$ is to the right of the zeros and $\rho_\varepsilon>0$, so that $Q_\varepsilon\ge0$ by~\eqref{dc:eq:Qw}, or some zero has $\beta=\frac12+\varepsilon$, and then $Q_\varepsilon(h)=\lim_{\varepsilon'\downarrow\varepsilon}Q_{\varepsilon'}(h)\ge0$, the limit holding because $Q_\varepsilon(h)$ is a finite combination of integrals continuous in $\varepsilon$.

Conversely, suppose $Q_\varepsilon\ge0$. The map $T_\varepsilon:g\mapsto\langle D,e^{-\varepsilon|x|}g\rangle$ is a distribution of order at most one, the corner of $e^{-\varepsilon|x|}$ at $0$ being harmless in the archimedean integral, and $Q_\varepsilon(h)=T_\varepsilon(h*\tilde h)\ge0$ says that it is positive definite. By the Bochner Schwartz theorem~\cite[Chapter II]{GelfandVilenkin4} it is tempered and $T_\varepsilon(g)=\int\hat g\,d\nu_\varepsilon$ for a positive tempered measure $\nu_\varepsilon$ and every Schwartz $g$. Fix $\varepsilon_0>\varepsilon$ with $\varepsilon_0\ge\frac12$, put $c=\varepsilon_0-\varepsilon$, and let $k_\eta(x)=e^{-c\sqrt{x^2+\eta^2}}$, smooth, positive definite and of exponential decay, with $\hat k_\eta\ge0$, $\hat k_\eta\to2\pi P_c$ pointwise and $\hat k_\eta\le C(1+\theta^2)^{-1}$ uniformly in $\eta\le1$. For $h\in C_c^\infty$ and $g=h*\tilde h$, $T_\varepsilon(k_\eta g)=\frac1{2\pi}\int(\hat k_\eta*|\hat h|^2)\,d\nu_\varepsilon$, and as $\eta\to0$ the left side tends to $T_\varepsilon(e^{-c|x|}g)=Q_{\varepsilon_0}(h)=\int|\hat h|^2\rho_{\varepsilon_0}$, because $k_\eta g\to e^{-c|x|}g$ uniformly, with supports in a fixed interval and uniformly bounded Lipschitz constants, which is enough for dominated convergence in the archimedean integral and in the other terms of $\langle D,\cdot\rangle$. By Fatou, $\int(P_c*|\hat h|^2)\,d\nu_\varepsilon\le\int|\hat h|^2\rho_{\varepsilon_0}<\infty$; choosing $h$ with $|\hat h|^2>0$ near $0$ gives $P_c*|\hat h|^2\ge c'(1+\theta^2)^{-1}$, so $\int(1+\theta^2)^{-1}d\nu_\varepsilon<\infty$, and dominated convergence then gives $\int(P_c*|\hat h|^2)\,d\nu_\varepsilon=\int|\hat h|^2\rho_{\varepsilon_0}$, that is $\int|\hat h|^2\,(P_c*\nu_\varepsilon)=\int|\hat h|^2\rho_{\varepsilon_0}$. Polarising, the same holds with $|\hat h|^2$ replaced by $\hat h_1\overline{\hat h_2}$, and these products span a dense set, so $\rho_{\varepsilon_0}=P_{\varepsilon_0-\varepsilon}*\nu_\varepsilon$. Theorem~\ref{tf:thm:deconv} gives $\beta\le\frac12+\varepsilon$ for every zero.
\end{proof}

The theorem is the Fourier dual of the deconvolution theorem; at $\varepsilon=0$ it is a proof of Weil's criterion along the lines of the harmonic continuation of Theorem~\ref{tf:thm:deconv}. The damping also explains why $\varepsilon=\frac12$ is free. The primes enter with weight $\Lambda(n)n^{-1/2-\varepsilon}$, absolutely summable exactly for $\varepsilon>\frac12$: the damped criterion at $\varepsilon=\frac12$ sits at the edge of the Euler product, the basepoint one of Chapter~\ref{chap:basepoint}.
